\chapter{层}

拓扑空间$X \in\textbf{Top}$上构造交换环层$\mathcal O_X$，
点$x \in X$上有局部环$\mathcal O_x$，
$\exists! \mathfrak m_x$，
构成局部赋环空间$(X,\mathcal O_X)$

实拓扑流形$X$上构造实连续函数层$\mathcal O_X$，
剩余域为实域．
实光滑流形$X$上构造实光滑函数层$\mathcal O_X$，
剩余域为实域．
复流形$X$上构造全纯函数层$\mathcal O_X$，
剩余域为复域．
概形$X$上构造正则函数层$\mathcal O_X$．

层范畴$\textbf{Sh}(X)$是Grothendieck topos，
$\mathcal O_X$视为其中的环对象．


\section{完备格}

\begin{framed}

$(L,\leq) \in\textbf{Poset}$是偏序集．
若对于$\forall a,b \in L$，
二元上确界(并join)$a\vee b$
和下确界(交meet)$a\wedge b$都存在，
即$\vee,\wedge$运算封闭，
称$L$为格．

若存在$0,1 \in L$ s.t.
$\forall x \in L$有$0 \leq x \leq 1$，
即$x \vee 0 = x,\ x \wedge 1 = x$，
或者等价的$x \vee 1 = 1,\ x \wedge 0 = 0$，
称$L$为有界格，
$0$为下确界，
$1$为上确界．

若任意子集$S \subseteq L$都有上确界
$\text{sup} \ S = \bigvee S$
和下确界
$\text{inf} \ S = \bigwedge S$，
称$L$为完备格．

\end{framed}

\begin{example}{$PX$}

幂集$PX$按照包含关系构成有界格和完备格，
格的交和并运算对应集合的交和并运算，
格有下确界$\varnothing$，
上确界$X$．

\end{example}

\begin{example}{$[0,1]$}

闭区间$[0,1]$按照自然序构成有界格和完备格，
格的交和并运算对应数的下确界和上确界运算．
格有下确界$0$，
上确界$1$．

\end{example}

\begin{example}{$\mathbb Z$}

$\mathbb Z$按照整除性$a \mid b$构成格，
$a \vee b = \text{lcm}(a,b)$，
$a \wedge b = \text{gcd}(a,b)$．
整数的无限子集，
例如素数的全体，
没有最小公倍数，
故$\mathbb Z$不是有界格，
更不是完备格．

若在理想上而非集合上考虑，
令$\mathcal I_\mathbb Z$为$\mathbb Z$的理想的全体．
有$a \mid b$ iff. $(b) \subseteq (a)$，
即便全体素数构成的无限子集$\{p_i\}$没有最小公倍数，
其理想$(p_i)$的交为$(0) \in \mathcal I_\mathbb Z$，
故下确界始终存在．
故$\mathcal I_\mathbb Z$是完备格．

\end{example}

交换环$A\in\textbf{CRing}$的理想全体
$\mathcal I_A$是有界格也是完备格．
以上例子告诉我们，
有时候在理想构成的完备格上，
相比环子集构成的格具有更好的性质．

\begin{framed}
格$(L,\vee,\wedge)$若满足以下相互等价的等式：
左分配律$a\wedge(b\vee c) = (a\wedge b) \vee (a\wedge c)$，
右分配律$a\vee(b\wedge c) = (a\vee b) \wedge (a\vee c)$，
则具有分配律．
\end{framed}

最常见的例子即是子集运算的分配律．

\section{Heyting代数}

\begin{framed}
有界分配格$(H,\vee,\wedge,0,1)$，
若有称之为蕴涵的二元封闭运算$\Rightarrow$，
$\forall a,b \in H$，
$\exists (a \Rightarrow b) \in H$
s.t. 
$\forall c \in H$，

$$
c \leq (a \Rightarrow b)
\Longleftrightarrow
(a \wedge c) \leq b
$$

称$(H,\vee,\wedge,\Rightarrow,0,1)$为Heyting代数．

定义补运算为：

$$
\neg a = (a \Rightarrow 0)
$$

\end{framed}

以上定义中的条件写为偏序范畴的态射集
$H(c,a \Rightarrow b) \simeq H(a \wedge c,b)$，
以便发现这里实际上是一对伴随函子$(a \wedge -) \dashv (a \Rightarrow -)$：

$$
% https://tikzcd.yichuanshen.de/#N4Igdg9gJgpgziAXAbVABwnAlgFyxMJZABgBoAmAXVJADcBDAGwFcYkQAKACVIB1fGMAI4BKEAF9S6TLnyEUZYtTpNW7bnwHCxk6djwEiAFgrKGLNohAAjCVJAZ9c46SU1zaq-QAE-AO4wUADmMN4AxnZ6soYo5K5mqpacPPyCopEOMgbyyHFU7onqKVrpuplOMcgAbPEFFuwRZY7ROTX5KvVevrwASlhBABY49ABOIxB+3rbiyoEhCCigAGbjALZIAIw0OBBIZCBwA1hLOHt1niCpwhkrEOuIAMzbu4hxHRdXQiA0jPTWMIwAApZZxWEb9IY3NZIEwgHZIACs5ySnyhdxhz02yPYPn8c1CAFpviBfv8gSCYiBwYNTmVbvd9vDEEj3klcb0IcMxhNvESfn8AcCKvIqZy0fcanCXk9WTjugFgoTiaTBRSRdTIXToYhJUyAOzYqyorXo16YxAG2Vdfh9GmjcaTPkkgXk4XsDW0yjiIA
\begin{tikzcd}
{(H,\leq)} \arrow[dd, "\leq"]             &  & {(H,\leq)} \arrow[dd, "\leq"] \arrow[ll, "a \wedge -"'] &  & a \wedge c \arrow[dd, "\leq"']   &  & c \arrow[ll, "a \wedge -"'] \arrow[dd, "\leq"] \\
                                          &  &                                                         &  &                                  &  &                                                \\
{(H,\leq)} \arrow[rr, "a \Rightarrow -"'] &  & {(H,\leq)}                                              &  & b \arrow[rr, "a \Rightarrow -"'] &  & a \Rightarrow b                               
\end{tikzcd}
$$

以上的蕴涵运算$a \Rightarrow b$常令人疑惑．
在经典逻辑中，
蕴涵写为$A \to B$，
直觉上理解为：
“若A真，则B真”．
这个陈述为假则意味着：
“A真而B假”即$A \wedge \neg B$，
除此之外都应该视为真：

$$
\neg (A \wedge \neg B) = \neg A \vee B
$$

这里的$\rightarrow$和$\neg$都符合Heyting代数的定义，
为防止二者的循环引用，
可以直接用真值表来定义蕴涵．

\begin{example}{$PX$}

幂集$PX$上，
有补集运算，
以及参照逻辑的方式定义蕴涵
$(A \Rightarrow -) = A^c \cup -$，
有
$C \subseteq A^C \cup B 
\Longleftrightarrow 
A \cap C \subseteq B$．

\end{example}

\begin{example}{$\text{Open} \ X$}

拓扑空间$X$上的开集范畴$\text{Open} \ X$
是$X$的子集族．
和$PX$对补的定义不同，
开集的补集不一定是开集，
为保持补运算的封闭性，
定义开集$A$的补$A^c$为$X \setminus A$中最大的开集，
再按照$PX$相似的方式定义蕴涵
$(A \Rightarrow -) = A^c \cup -$，

$C \subseteq A^C \cup B 
\Longleftrightarrow 
A \cap C \subseteq B$．

\end{example}

Heyting代数上自带的伴随函子，
使得它构成CCC．
蕴涵$\Rightarrow$可以构造指数对象．


\section{生成理想}

理想作为子集可以覆盖给定的子集$S$，
而理想又可以被素理想所覆盖，
这样就引入了子集$S$所生成的理想，
以及相应的素理想问题．

\begin{example}{$\mathbb Z$}

子集$S=\{ -60, 120, 720 \}$，
覆盖它的理想集合为$\{(1),(2),(3),(5),(6),(15),(60)\} \in I_\mathbb Z$，
最大公约数$\text{gcd} \ S = 60$，
理想$(\text{gcd}) = (60) = (1)\cap(2)\cap(3)\cap(5)\cap(6)\cap(15)\cap(60)$．

$$
\begin{tikzcd}
     &  &                  &  & (2) \arrow[lld] \arrow[lllldd, bend right] \\
     &  & (6) \arrow[lld]  &  &                                            \\
(60) &  &                  &  & (3) \arrow[llu] \arrow[lld] \arrow[llll]   \\
     &  & (15) \arrow[llu] &  &                                            \\
     &  &                  &  & (5) \arrow[llu] \arrow[lllluu, bend left] 
\end{tikzcd}
$$

按照理想升链，
跳过中间过程的理想$(6),(15)$，
$S$最终被极大素理想$(2),(3),(5)$所覆盖，
记

$$
V(S) = \{(2),(3),(5)\}
$$

取整系数$\lambda_i\in \mathbb Z$做有限形式和
$\lambda_1 \cdot (-60) + \lambda_2 \cdot 120 + \lambda_3 \cdot 720$，
其全体构成理想$(S) = \{ \lambda_1 \cdot (-60) + \lambda_2 \cdot 120 + \lambda_3 \cdot 720 \}$．

\end{example}

\begin{example}{$\mathbb R[X]$}

$X = \mathbb R$，
多项式环$\mathbb R[X]$的子集
$S=\{ f_1(x) = x - a, f_2(x) = x - b, f_3(x) = x-3 \}$，
有零点$Z(S) = \{ a, b, c\}$．
$S$的任何子集$S_i$的零点集
$Z(S_i)$都覆盖了$S$，
这些零点集的交为：
$\bigcap\limits_i Z(S_i)$

$$
% https://tikzcd.yichuanshen.de/#N4Igdg9gJgpgziAXAbVABwnAlgFyxMJZAFgBoAGAXVJADcBDAGwFcYkQAtACnoEoQAvqXSZc+QigBMpAIzU6TVu271SAI35CR2PASLlSk+QxZtEnHutIBjTcJAYd4omSM0TS89w2D7jsXpSpADMxopmFmo2dtoBEiSkxGGmyly2gvIwUADm8ESgAGYAThAAtkgGIDgQSDLu4UhgzIyMNIz0ajCMAAqiuhIgjDAFOL6FJeWIldVI0gopiE0tbR1dvU6Bg8OjNJ1gUEgAtMHkWiDFZbU0M4hzHmZLrYOrPX3O5kMjY+cTSMHXNUQdXmnkeK06rw2A0+ozOF0m-yqgLIIIezSe7Qh6zi7Bh33hfwBs3qCzBzyxb02eLhv0QKJudwai3R4LWlOh23xtIArES6STQSzyWyobjOTTLoheUjiajGkLMSKcR9ObsYPs-qdKAIgA
\begin{tikzcd}
           &  &                      &  & Z(a) \arrow[lld] \arrow[lllldd, bend right] \\
           &  & {Z(a,b)} \arrow[lld] &  &                                             \\
{Z(a,b,c)} &  &                      &  & Z(b) \arrow[llu] \arrow[lld] \arrow[llll]   \\
           &  & {Z(b,c)} \arrow[llu] &  &                                             \\
           &  &                      &  & Z(c) \arrow[llu] \arrow[lllluu, bend left] 
\end{tikzcd}
$$

取环系数$\forall \lambda_i \in A$
做有限和$\sum_i \lambda_i f_i \in A$，
其全体构成理想
$(S) = \{ \sum_i \lambda_i f_i \in A \}$．

\end{example}

子集$S$生成的理想$S$有如下两种定义：

\begin{framed}
\noindent
$(S) = \bigcap\limits_{I \supset S,\ I \in I_A} I$
\end{framed}

\begin{framed}
\noindent
$(S) = \left\{ \sum_{i=1}^n a_is_i \mid n \in \mathbb N,\ (a_i,s_i) \in A \times S \right\}$
\end{framed}

这样可以引入主理想和PID的概念：

\begin{framed}
\noindent
由单个元素 $a \in A \in \textbf{CRing}$ 生成的理想 $(a) = \{ra \mid r \in A\}$ 称为\textbf{主理想}（Principal Ideal）．
若整环的所有理想都是主理想，则称其为\textbf{主理想整环}（PID）．
\end{framed}

$\mathbb{Z}$ 是 PID 的典型例子；而多项式环 $\mathbb{Z}[x]$ 并非 PID（例如理想 $(2, x)$ 不能由单元素生成）．

\begin{framed}
\noindent
由单个元素 $a \in A \in \textbf{CRing}$ 生成的理想 $(a) = \{ra \mid r \in A\}$ 称为\textbf{主理想}（Principal Ideal）．
若整环的所有理想都是主理想，则称其为\textbf{主理想整环}（PID）．
\end{framed}

\section{反单调Galois连接}

$A$的幂集有自然的偏序结构$(\textbf 2^A,\leq)$，
把子集$S \in \textbf 2^A$置于这个偏序结构考虑．
对函数环$A = C(X)$的子集$S$求零点集，
视为幂集上的函数：
$Z: \textbf 2^A \to \textbf 2^X:: S \mapsto Z(S)$．
反过来，
子集$E \in X$确定了理想
$I(E) = \{f \in A \mid \forall x \in E:\ f(x) = 0\}$，
视为幂集上的函数：
$I: \textbf 2^X \to \textbf 2^A :: E \mapsto I(E)$．

前面讨论过，
$\textbf 2^X,\textbf 2^A$是偏序集和完备格，
但我们需要讨论的是$A$中理想之间的结构，
而非$A$的子集之间的结构，
需要转换到$I_A$这个完备格上讨论．
将子集$S\in\textbf 2^A$用$(S)\in \text{Ideal} \ A$来替换，
得到
$Z: I_A \to \textbf 2^X:: (S) \mapsto Z((S)) = Z(S)$，
以及
$I: \textbf 2^X \to I_A :: E \mapsto I(E)$．

\begin{framed}
\noindent

反单调Galois连接是连接两个偏序范畴$(P,\leq)$和$(Q,\leq)$的反变的伴随函子，
$q \leq Lp$ iff. $p \leq Rp$对应了
$(Q,\leq)(q,Lp) \simeq (P,\leq)(p,Rq)$：

$$
% https://tikzcd.yichuanshen.de/#N4Igdg9gJgpgziAXAbVABwnAlgFyxMJZABgBoAmAXVJADcBDAGwFcYkQAKARVIB1fGMAI4BKAHr8cMAB45gENAF8Qi0uky58hFGWLU6TVu258Bw8ZJlyFy1eux4CRACwV9DFm0QghKtSAwHLRdSPRoPI28AGTQ-e00nFHJQ90MvTgAFU0FROICNR21kZKpwtOMs-hyRPMCEooA2FLLPdli7fKDE5CbSg1bvACVfRX0YKABzeCJQADMAJwgAWyQyEBwIJABGGjgACyxZnFWWyJAq4TyF5aRk9c3EAGZT9IvfGkZ6ACMYRgyC4LeQRHK6LFaIVz3JAAVhe7DeoJuEJoG22cOiIA+31+-y62hA8ywEz2xw613Ba1RiFh-TOg0xIE+Pz+AMSBKJJMR4KaUKe6JAUQZTJxrPxhOJpP85KQPKpAHZ+QiyWDbiiHgrael6Vjmbj6uxxZzRoogA
\begin{tikzcd}
{(Q,\leq)^\text{op}}                                     &  & {(P,\leq)} \arrow[dd, "\leq"] \arrow[ll, "L"'] &  & Lp                                    &  & p \arrow[ll, "L"'] \arrow[dd, "\leq"] \\
                                                         &  &                                                &  &                                       &  &                                       \\
{(Q,\leq)^\text{op}} \arrow[uu, "\leq"] \arrow[rr, "R"'] &  & {(P,\leq)}                                     &  & q \arrow[uu, "\leq"] \arrow[rr, "R"'] &  & Rq                                   
\end{tikzcd}
$$

\end{framed}

这样$(Z,I)$构成反单调Galois连接：

$$
% https://tikzcd.yichuanshen.de/#N4Igdg9gJgpgziAXAbVABwnAlgFyxMJZABgBoAmAXVJADcBDAGwFcYkQAKAHS5xgA8cAAnIA9ABqkejGAEcAlKJ59BwCGgC+IDaXSZc+QijLFqdJq3bdeA4WMnS5i5bbWbtukBmx4CRACwUZgwsbIggAKIeej6GAaSmNCGW4QBaHBwAyvLy0V76vkbI5AnBFmGcAJIA+gCCUlwyCnneBn4oJVRJ5VY19Y7NOjFtRQBspd2hVtktBXEo413mU+GVHBG5GmYwUADm8ESgAGYAThAAtkhkIDgQSACMNHAAFlhHOFeTKSADeacXSBKNzuiAAzF8Kr8aIx6AAjGCMAAKc3aIBk7z+Z0uiECwKQAFYIexfkMQP9sbjbg8iWkQNC4QjkbFUScsLtnh9SeTPnjEITlt9KnS0QykSijCBWezOZ5uYhxrzwQKKqlhTD4WLmRKpRzMQD5TQqYgAOw0n6NOR67FAo2m5XsIX0jVMkbsHWcygaIA
\begin{tikzcd}
{(\text 2^X,\leq)^\text{op}}                                     &  & {(I_A,\leq)} \arrow[dd, "\leq"] \arrow[ll, "Z"'] &  & Z((S))                                &  & (S) \arrow[ll, "Z"'] \arrow[dd, "\leq"] \\
                                                                 &  &                                                  &  &                                       &  &                                         \\
{(\text 2^X,\leq)^\text{op}} \arrow[uu, "\leq"] \arrow[rr, "I"'] &  & {(I_A,\leq)}                                     &  & E \arrow[uu, "\leq"] \arrow[rr, "I"'] &  & I(E)                                   
\end{tikzcd}
$$

\section{函数环}

实微分几何的基本研究对象是实光滑流形$X$上的实光滑函数环$C(X)$．
在$X$上选定子集$S$得到零点集 
$Z: \textbf 2^X \to \textbf 2^{C(X)} :: S \mapsto Z(S) = \{ f \in C(X) \mid \forall x \in S,\ f(x) = 0 \}$．
这构造了幂集之间的反单调偏序Galois连接．
除非$X$是单点，否则$C(X)$不是整环．
    
每个点$x \in X$都可以构造退化函数集合作为极大理想
$X \to \text{Max} \ C(X) : x \mapsto \mathfrak m_x = Z(\{ x \}) = \{ f \in C(X) \mid f(x) = 0 \}$．
构造$x$点上的赋值同态：

$$
\begin{aligned}
\varphi: C(X)/\mathfrak m_x &\to \mathbb R \\
f + \mathfrak m_x &\mapsto \varphi(f + \mathfrak m_x) = f(x) \\
\end{aligned}
$$

常值函数可以取值任意实数，
故$\varphi$是满的．
同态的核$\text{Ker} \ \varphi = \{ 0 \}$，
故$\varphi$是单的，
故$\varphi$是同构即

\begin{equation}
C(X)/\mathfrak m_x \simeq \mathbb R
\end{equation}
\label{eq:fun-ring-residue-field-real}

进而构造乘法系统$C(X)\setminus \mathfrak m_x$，
局部化：
$\mathcal O_x = C(X)_{\mathfrak{m}_x}$，
其元素为分式 $f/g$，
其中 $(f,g \in C(X),\ g(x) \neq 0$．
$f/g = f'/g'$ iff. $\exists h \neq 0$ s.t. $h(f g' - f' g) = 0$．
即在点 $x$ 的某个邻域上 $f/g = f'/g'$．
因此 $\mathcal O_x$ 恰好同构于 $x$ 点的光滑函数芽环，
即所有在 $x$ 的某个邻域上有定义的光滑函数在“局部相等”关系下的等价类构成的环．
根据(\ref{eq:fun-ring-residue-field-real})，
类似(\ref{eq:short-exact-seq-num-field})：

\begin{equation}
0 \to \mathfrak m \to \big( \mathcal O_x = C(X)_{\mathfrak{m}_x} \big) \to \mathbb R \to 0
\end{equation}
\label{eq:short-exact-seq-fun-real}


\section{函数层}

拓扑空间$X$的开集族$\text{Open} \ X$在嵌入包含下构成偏序范畴，
函数层的构造是反变的偏序函子$\mathcal O$：

$$
% https://tikzcd.yichuanshen.de/#N4Igdg9gJgpgziAXAbVABwnAlgFyxMJZABgBpiBdUkANwEMAbAVxiRAAoAdTnGADxzAA8mhhgAvgAJukgBqlpnOEwBGcGLwCOASgB63XgOAQ04kONLpMufIRQAmclVqMWbLj344VAM2ABhACUsMABzC0VDQQAneHFtc0sQDGw8AiIyAEZnemZWRBAAVUSrVNsiR2zqXLcC7gBbOhwACwBjRkkhAH1CyQBeSX9ii1KbdJQyAGYc13yQADUS5Os0u2RHaerZtgamto7u+f7BxfFnGChQ+CJQH2iIeqQyEBwIJEytvJ3ORpb2hk6XVkSzuDyQjheb0QkxGIFBj0QABZqK8kABWWHw8EoqHIlxfOpKVTqLQgagMOgqGAMAAKK3KBWiWFCzRwIPuCLROKQk0+tRABi8wFicDM5Mp1LpZXGICZLLZZ3EQA
\begin{tikzcd}
{(\text{Open} \ X, \subseteq)^\text{op}} \arrow[rr, "\mathcal O_X"] &  & {(\textbf{CRing}, \text{res})}              \\
U \arrow[rr] \arrow[dd, "\subseteq"']                               &  & \mathcal O_U = CU                           \\
                                                                    &  &                                             \\
V \arrow[rr]                                                        &  & \mathcal O_V = CV \arrow[uu, "\text{res}"']
\end{tikzcd}
$$

限制映射$\text{res}: \mathcal O_V \to \mathcal O_U :: f \mapsto f\mid_U$．
函子$\mathcal O_X$是$X$上的 \textbf{结构层} (structure sheaf)．
这样的函子的全体构成函数层范畴$\textbf{Sh}(\textbf{CRing})$．

为了考虑$X$上的局部性质，
需要让局部以外的性质被忽略等价商掉．
比较函数$f,g\in\mathcal O_X = C(X)$，
其差异$f-g$的全局性质，
随着限制的极限过程而消去：
$(f-g)\mid_{U_1} \stackrel{\text{res}}{\longrightarrow}(f-g)\mid_{U_2} \stackrel{\text{res}}{\longrightarrow} \cdots $，
在这一极限过程中，
若存在点$x$的开邻域$U$ s.t. $(f-g)\mid_U = 0$，
这意味着$f,g$在$x$点的局部是相等的，
构成等价关系$f \sim_x g$，
等价类称为函数芽$[f]_x = [g]_x$，
芽构成的环称为\textbf{茎}(stalk)：
$\mathcal O_x = \varinjlim_{U \ni p} \mathcal O_U$．

在$x$处退化的函数芽构成极大理想
$\mathfrak m_x = \{ [f] \in \mathcal{O}_x \mid f(x) = 0 \}$．
由于茎的构造基于拓扑空间$X\in\textbf{Top}$，
函数的连续性是基本要求，
若不加说明函数环都是连续函数环．
在极大理想之外的芽$f \notin \mathfrak m_x$是非退化的，
由于函数的连续性保证了$f$的可逆性．

这是唯一的极大理想，
使得茎$\mathcal O_x$是局部环，
以及 $(X, \mathcal O_X)$ 构成一个\textbf{局部赋环空间}（locally ringed space）．


